% ============================================
% MEM-03 (ZH)
% 中文翻译 — SPEC-935-R356
% ============================================
\section*{\textbf{MEM-03} --- 父亲的死}
\markboth{MEM-03 --- 父亲的死}{}

\noindent \textit{没有水的第二天}

\noindent \textbf{1915 —— 没有水的第二天}

\vspace{0.5em}\noindent\rule{\textwidth}{0.5pt}\vspace{0.5em}

\noindent 父亲在一棵胡亚泽罗树下坐了下来。

\noindent 树已经枯了——没有叶子，没有树荫。但他还是坐下了。背靠着树干，闭上眼睛。

\noindent \textit{“爹，走吧。”}

\noindent \textit{“你们走吧。”}

\noindent 母亲跪到他面前。

\noindent \textit{“安东尼奥，起来。”}

\noindent \textit{“起不来了。”}

\noindent \textit{“起得来。你总说起得来。”}

\noindent 他睁开眼睛。看着她。他的目光很远。仿佛他已经到了别的地方。

\noindent \textit{“玛丽亚，我感觉不到我的腿了。”}

\noindent 她看着他的腿。腿肿着。脚已经裂开，脚后跟上有干涸的血。

\noindent \textit{“那是累的，”她说。}

\noindent \textit{“那是死。”}

\noindent \textit{“别这么说。”}

\noindent \textit{“玛丽亚。”他把手放在她的脸上。他的手只剩骨头。“带上孩子。走吧。别停。”}

\noindent \textit{“我不会把你丢在这里。”}

\noindent \textit{“我已经死了。只差身体停下来。”}

\noindent 母亲开始哭。那是我第一次看见她哭。她的哭声是干的——没有眼泪。身体里没有水可以做成眼泪。

\noindent \textit{“若阿金，”父亲说，“过来。”}

\noindent 我走过去。

\noindent 他把手放在我的头上。

\noindent \textit{“你现在是男子汉了。你要照顾你娘。”}

\noindent \textit{“是，爹。”}

\noindent \textit{“你不能让她挨饿。”}

\noindent \textit{“……是，爹。”}

\noindent \textit{“该吃什么就吃什么。哪怕……”他停住了。望着地平线。“哪怕那东西看起来不像食物。”}

\noindent \textit{“是，爹。”}

\noindent 他把我的脸拉到他的脸前。他的气息闻起来像空胃。

\noindent \textit{“你不会饿死，若阿金。你听见了吗？你不会饿死。该做什么就做什么。但你不能饿死。”}

\noindent \textit{“是，爹。”}

\noindent 他松开我的脸。闭上眼睛。

\noindent \textit{“可以走了。”}

\noindent 母亲站着没动。

\noindent \textit{“可以走了，”他重复道，“我说了。”}

\noindent 她站起来。抓住我的手。

\noindent 我们走了。

\noindent 当我回头时，他还坐在那棵枯干的胡亚泽罗树下。头垂着。手臂伸在膝上。

\noindent 母亲拉了拉我的手。\textit{“走吧。”}我本可以转过脸去。她没有回头看。但我选择把那个画面留下来。最后一个。

\noindent 我最后一次看见父亲时，他像一位圣人。

\noindent 一位饿死的圣人。

\vspace{0.5em}
\noindent\textit{\textbf{↪ 路线:} \rota{MEM-02} \rotasep \rota{MEM-07} \rotasep \rota{MEM-04}}
